% load packages
\usepackage{geometry}
\usepackage{xcolor}
\usepackage{eso-pic}
\usepackage{fancyhdr}
\usepackage{sectsty}
\usepackage{fontspec}
\usepackage{titlesec}
\usepackage{graphicx}
\usepackage{tikz}
\usepackage{amsmath}
\usepackage{amssymb}
\usepackage{unicode-math}
  
%% Set page size with a wider right margin
\geometry{a4paper, total={170mm,257mm}, left=25mm, top=20mm, bottom=20mm, right=25mm}

%% Let's define some colours
\definecolor{light}{HTML}{F4F6F8}   % code-block background (_extension.yml: code-block-bg: light)
\definecolor{highlight}{HTML}{4F5B66}   % links (dark slate grey)
\definecolor{dark}{HTML}{000000}
\definecolor{acecyan}{HTML}{03E5FF}
\definecolor{acemagenta}{HTML}{E262E2}
\definecolor{acegrey}{HTML}{5A6570}   % caption/note grey, same as the site

% ---- Tables (tinytable / tabularray) in the site's look ---------------------
% Cyan top rule, light grey header and bottom rules (table notes are styled
% further down, with the captions). tabularray refuses to redefine a table command, so the booktabs rules
% are recolored by replacing their implementations; every step is guarded, so
% a tabularray version with other internals just keeps its default rules.
\definecolor{acerule}{HTML}{BFC5CA}
\ExplSyntaxOn
\AtBeginDocument{
  \bool_lazy_and:nnT
    { \cs_if_exist_p:c { __tblr_table_command_toprule:w } }
    { \cs_if_exist_p:c { tblr@booktabs@hline } }
    {
      \RenewDocumentCommand \__tblr_table_command_toprule:w { O{} }
        { \use:c { tblr@booktabs@hline } [wd=0.12em, fg=acecyan, #1] }
      \RenewDocumentCommand \__tblr_table_command_midrule:w { O{} }
        { \use:c { tblr@booktabs@hline } [wd=0.05em, fg=acerule, #1] }
      \RenewDocumentCommand \__tblr_table_command_bottomrule:w { O{} }
        { \use:c { tblr@booktabs@hline } [wd=0.05em, fg=acerule, #1] }
    }
  \bool_lazy_and:nnT
    { \cs_if_exist_p:c { __tblr_table_command_cmidrule:w } }
    { \cs_if_exist_p:c { tblr@booktabs@cline } }
    {
      \RenewDocumentCommand \__tblr_table_command_cmidrule:w { O{} m }
        { \use:c { tblr@booktabs@cline } [wd=\cmidrulewidth, endpos, fg=acerule, #1] {#2} }
    }
}
\ExplSyntaxOff

% ---- Markdown (pandoc longtable) tables in the same look --------------------
% booktabs colours its rules with \CT@arc@ (colortbl's hook, empty without
% colortbl), so each rule sets it first: a 0.12em cyan top rule, 0.05em light
% grey rules otherwise (pandoc never passes these rules a width of its own). tabularray's own copies of these commands are handled above.
% \aceTableHead is inserted by _filters/ace-tables.lua.
\definecolor{acelabel}{HTML}{767676}
\definecolor{acehair}{HTML}{E6E9EC}
% Table type sizes, the same for both kinds of table: body at 0.9em and column
% labels at 0.85em of the normal size (what ace_tt_pdf() and yesno_tt() give
% tinytable), notes and captions \small (set below with the captions).
\newcommand{\aceTableHead}{\normalsize\ttfamily\fontsize{0.85em}{1.1em}\selectfont\color{acelabel}}
\makeatletter
\AtBeginDocument{%
  \@ifundefined{@BTrule}{}{%
    \providecommand*\CT@arc@{}%
    \let\aceBT@toprule\toprule
    \let\aceBT@midrule\midrule
    \let\aceBT@bottomrule\bottomrule
    \let\aceBT@cmidrule\cmidrule
    \def\toprule{\noalign{\gdef\CT@arc@{\color{acecyan}}}\aceBT@toprule[0.12em]}%
    \def\midrule{\noalign{\gdef\CT@arc@{\color{acerule}}}\aceBT@midrule[0.05em]}%
    \def\bottomrule{\noalign{\gdef\CT@arc@{\color{acerule}}}\aceBT@bottomrule[0.05em]}%
    \def\cmidrule{\noalign{\gdef\CT@arc@{\color{acerule}}}\aceBT@cmidrule}%
  }%
}
\makeatother
% Body size and a little more air between rows of a Markdown table, to match tinytable.
\makeatletter
\@ifundefined{AddToHook}{}{\AddToHook{env/longtable/begin}{\fontsize{0.9em}{1.2em}\selectfont\renewcommand{\arraystretch}{1.3}}}
\makeatother


% Facet badges on the cover, matching the chips on the website: the fill is
% the site's translucent tint flattened onto white, the line is the facet
% hue, the ink is the darker text tone used for the facet key.
% The badges live on the cover, which is dark, so these are the site's dark
% theme values: the tint flattened onto the cover's background, the hue
% itself for the outline, and the brighter ink the site uses on dark.
\definecolor{acefacetpopfill}{HTML}{0D3941}\definecolor{acefacetpopline}{HTML}{03E5FF}\definecolor{acefacetpopink}{HTML}{03E5FF}
\definecolor{acefacettrtfill}{HTML}{1A313F}\definecolor{acefacettrtline}{HTML}{4DB9F5}\definecolor{acefacettrtink}{HTML}{6CC4FF}
\definecolor{acefacetoutfill}{HTML}{28293D}\definecolor{acefacetoutline}{HTML}{978EEC}\definecolor{acefacetoutink}{HTML}{A99EF2}
\definecolor{acefacetdesfill}{HTML}{35213C}\definecolor{acefacetdesline}{HTML}{E262E2}\definecolor{acefacetdesink}{HTML}{E39BE0}
\definecolor{acefacettext}{HTML}{F2F2F2}   % term text, as on the site in dark
\definecolor{aceink}{HTML}{111418}     % body text on the light interior pages

% The cover is dark, and carries its own small palette. Everything on it is
% drawn in these colors, so the interior pages are untouched by them.
\definecolor{acecoverbg}{HTML}{0F1317}    % the page itself
\definecolor{acecoverink}{HTML}{F4F6F7}   % titles, names, values
\definecolor{acecovermute}{HTML}{9AA4AE}  % keys, affiliations, notes
\definecolor{acecoversub}{HTML}{C7D1D9}   % the subtitle, a step brighter
\definecolor{acecoverlink}{HTML}{6CC4FF}  % links
% Band behind the cover title: a shade lifted off the page, not a panel.
\definecolor{acetitleband}{HTML}{1B2127}
\definecolor{aceecg}{HTML}{2C2F33}      % cover ECG trace, a step above the band
% Mono caps, letterspaced, as the chips and the sidebar facet titles are.
\newcommand{\aceFacetKeyFont}{\fontsize{9}{10.5}\selectfont\ttfamily\bfseries\addfontfeature{LetterSpace=14}}
\newcommand{\aceCoverLabelFont}{\fontsize{10}{12}\selectfont\ttfamily\bfseries\addfontfeature{LetterSpace=18}}
\newcommand{\aceCoverIDFont}{\fontsize{9.5}{11}\selectfont\ttfamily\addfontfeature{LetterSpace=12}}
\newcommand{\aceCoverDateFont}{\fontsize{11.5}{13.5}\selectfont\ttfamily\bfseries\addfontfeature{LetterSpace=4}}
\newcommand{\aceInfoValueFont}{\fontsize{9.5}{12}\selectfont}
\newcommand{\aceInfoMonoFont}{\fontsize{9.5}{11}\selectfont\ttfamily}
\newcommand{\aceInfoNoteFont}{\fontsize{8.5}{10.5}\selectfont}
\newcommand{\aceFacetTermFont}{\fontsize{9.4}{11}\selectfont\ttfamily\bfseries\addfontfeature{LetterSpace=5}}

% ---------------------------------------------------------
% CODE BLOCKS AND SYNTAX COLOURS
% Same palette as the site's light theme (see "CODE: inline and
% blocks" in styles.css): a very light panel, near-black text,
% and the ACE token colours.
% ---------------------------------------------------------
\definecolor{acecodebg}{HTML}{F4F6F8}    % panel
\definecolor{acecodefg}{HTML}{1F2933}    % normal text, variables
\definecolor{acecodefun}{HTML}{1A1F24}   % functions, builtins
\definecolor{acecodekw}{HTML}{8A2D8A}    % keywords, control flow, types
\definecolor{acecodestr}{HTML}{0B6B8F}   % strings
\definecolor{acecodenum}{HTML}{B0296F}   % numbers, constants
\definecolor{acecodecom}{HTML}{7B8590}   % comments
\definecolor{acecodeop}{HTML}{5B6670}    % operators, attributes, other
\definecolor{acecodeerr}{HTML}{C92A2A}   % alerts and errors

\makeatletter
% The token macros exist only when a document contains highlighted code, so
% every override is guarded and applied at the start of the document, after
% pandoc has defined them.
\newcommand{\aceTokPlain}[1]{#1}
\newcommand{\aceSetTok}[3]{%
  \@ifundefined{#1}{}{%
    \expandafter\renewcommand\csname #1\endcsname[1]{\textcolor{#2}{#3{##1}}}}}
\AtBeginDocument{%
  \definecolor{shadecolor}{HTML}{F4F6F8}%
  \definecolor{codeblockbg}{HTML}{F4F6F8}%
  \@ifundefined{fvset}{}{\fvset{fontsize=\small}}%
  \aceSetTok{NormalTok}{acecodefg}{\aceTokPlain}%
  \aceSetTok{VariableTok}{acecodefg}{\aceTokPlain}%
  \aceSetTok{FunctionTok}{acecodefun}{\textbf}%
  \aceSetTok{BuiltInTok}{acecodefun}{\textbf}%
  \aceSetTok{ExtensionTok}{acecodefun}{\textbf}%
  \aceSetTok{KeywordTok}{acecodekw}{\aceTokPlain}%
  \aceSetTok{ControlFlowTok}{acecodekw}{\aceTokPlain}%
  \aceSetTok{ImportTok}{acecodekw}{\aceTokPlain}%
  \aceSetTok{DataTypeTok}{acecodekw}{\aceTokPlain}%
  \aceSetTok{StringTok}{acecodestr}{\aceTokPlain}%
  \aceSetTok{VerbatimStringTok}{acecodestr}{\aceTokPlain}%
  \aceSetTok{SpecialStringTok}{acecodestr}{\aceTokPlain}%
  \aceSetTok{CharTok}{acecodestr}{\aceTokPlain}%
  \aceSetTok{DecValTok}{acecodenum}{\aceTokPlain}%
  \aceSetTok{FloatTok}{acecodenum}{\aceTokPlain}%
  \aceSetTok{BaseNTok}{acecodenum}{\aceTokPlain}%
  \aceSetTok{ConstantTok}{acecodenum}{\aceTokPlain}%
  \aceSetTok{CommentTok}{acecodecom}{\textit}%
  \aceSetTok{DocumentationTok}{acecodecom}{\textit}%
  \aceSetTok{AnnotationTok}{acecodecom}{\textit}%
  \aceSetTok{CommentVarTok}{acecodecom}{\textit}%
  \aceSetTok{InformationTok}{acecodecom}{\textit}%
  \aceSetTok{OperatorTok}{acecodeop}{\aceTokPlain}%
  \aceSetTok{OtherTok}{acecodeop}{\aceTokPlain}%
  \aceSetTok{SpecialCharTok}{acecodeop}{\aceTokPlain}%
  \aceSetTok{AttributeTok}{acecodeop}{\aceTokPlain}%
  \aceSetTok{PreprocessorTok}{acecodeop}{\aceTokPlain}%
  \aceSetTok{AlertTok}{acecodeerr}{\textbf}%
  \aceSetTok{ErrorTok}{acecodeerr}{\textbf}%
  \aceSetTok{WarningTok}{acecodenum}{\textit}%
}
\makeatother

%% Figure and table captions: small grey text, bold grey label.
\usepackage{caption}
\captionsetup{font={small,color=acegrey}, labelfont={bf,color=acegrey}, labelsep=period}

%% Table notes (tinytable renders them through tabularray): same grey.
\makeatletter
\AtBeginDocument{%
  \@ifpackageloaded{tabularray}{%
    \SetTblrStyle{note}{font=\small\color{acegrey}}%
    \SetTblrStyle{note-tag}{font=\small\color{acegrey}}%
    \SetTblrStyle{remark}{font=\small\color{acegrey}}%
    \SetTblrStyle{remark-tag}{font=\small\color{acegrey}}%
    \SetTblrStyle{caption}{font=\small\color{acegrey}}%
    \SetTblrStyle{caption-tag}{font=\small\bfseries\color{acegrey}}%
    \SetTblrStyle{caption-sep}{font=\small\color{acegrey}}%
  }{}%
}
\makeatother
%% Free-text notes under tables (citation_table() emits \begin{acetablenote}).
\newenvironment{acetablenote}{\par\small\color{acegrey}\setlength{\parindent}{0pt}\noindent\ignorespaces}{\par}

% Report ID placeholder, populated from YAML by fields.lua
\providecommand{\aceID}{}
\providecommand{\aceDOI}{}
\providecommand{\aceConceptDOI}{}
\providecommand{\coverAuthorsBlock}{}
\providecommand{\coverCategoriesBlock}{}
\providecommand{\aceAuthorsLine}{}
\providecommand{\aceAffiliationsLine}{}
\newlength{\aceInfoY}
\newlength{\aceInfoTop}
\newsavebox{\aceTitleBox}
\newlength{\aceTitleHeight}
\newlength{\aceTitleWidth}
\newlength{\aceTitleRun}
\newlength{\aceBandTop}
\newlength{\aceBandBottom}
\newlength{\aceTitleMeasure}
\setlength{\aceTitleMeasure}{16.6cm}

% One PQRST beat for the cover trace, drawn from the current point and 3.2cm
% wide: baseline, a short P wave, the QRS spike, then a short T wave back to
% baseline. P and T are single Bezier curves, so they read as smooth humps and
% stay well under the spike; only the QRS complex is drawn with straight lines,
% because that part of the signal really is angular. \aceEcgStrip chains seven
% beats into ONE path: drawn as separate paths, the round caps where two beats
% met overlapped into a visible bright dot on the baseline.
\newcommand{\aceEcgBeat}{%
  -- ++(0.66,0)
  .. controls +(0.12,0.08) and +(-0.12,0.08) .. ++(0.34,0)
  -- ++(0.38,0)
  -- ++(0.07,-0.10) -- ++(0.09,0.86) -- ++(0.10,-1.02) -- ++(0.09,0.26)
  -- ++(0.36,0)
  .. controls +(0.17,0.11) and +(-0.17,0.11) .. ++(0.50,0)
  -- ++(0.61,0)}
\newcommand{\aceEcgStrip}{%
  \aceEcgBeat\aceEcgBeat\aceEcgBeat\aceEcgBeat\aceEcgBeat\aceEcgBeat\aceEcgBeat}

% The glitch corner wedge and the hexagon logo that used to be painted into
% the page background were removed: the interior pages carry the mark in the
% running head instead (see PAGE FURNITURE below), which keeps them aligned
% with the text block and out of the way of tables and figures.


% ---------------------------------------------------------
% PAGE NUMBERING & HEADER STYLING
% ---------------------------------------------------------

% The ACE mark sits in the top corner of every page after the cover, with the
% report ID tagged beneath it in white on grey; the folio sits alone at the
% foot. brand/logo_mark is the wordmark trimmed to its own ink, so the width
% below is the width that prints. The block is \smash-ed, so its height never
% reaches the text block, and \aceRunDrop is tuned to leave the mark about
% 6 mm from the trim and the tag about 3 mm clear of the first body line.
\newcommand{\aceRunMarkW}{1.7cm}
\newlength{\aceRunDrop}\setlength{\aceRunDrop}{3mm}
\newlength{\aceRunOut}\setlength{\aceRunOut}{19mm}  % how far the block hangs into the right margin
\newcommand{\aceRunIDBox}{%
  \setlength{\fboxsep}{3pt}%
  \colorbox{acegrey}{\textcolor{white}{\ttfamily\scriptsize\aceID}}}
\newcommand{\aceRunBlock}{%
  \smash{\raisebox{-\aceRunDrop}{\shortstack{%
    \includegraphics[width=\aceRunMarkW]{brand/logo_mark.png}\\[3pt]
    \aceRunIDBox}}}}
% "Page x of N". lastpage records the final page number; \pageref* (the
% hyperref form) keeps the total from being drawn as a link.
\usepackage{lastpage}
% \needspace{<lines>} starts a new page when less room than that is left (keeps short tables whole).
\usepackage{needspace}
% No table or figure may float past the heading of the next section or
% subsection, so every float stays inside the (sub)section that refers to it.
\usepackage{placeins}
\usepackage{etoolbox}
\pretocmd{\section}{\FloatBarrier}{}{}
\pretocmd{\subsection}{\FloatBarrier}{}{}
\newcommand{\aceLastPage}{\ifdefined\hypersetup\pageref*{LastPage}\else\pageref{LastPage}\fi}
\newcommand{\aceFolio}{{\color{acegrey}\ttfamily\footnotesize Page \thepage\ of \aceLastPage}}

\setlength{\headheight}{26pt}
\setlength{\headsep}{14pt}
\setlength{\footskip}{18pt}   % folio ~12mm from the trim, two lines clear of the text

\fancypagestyle{mystyle}{
  \fancyhf{}
  \renewcommand\headrulewidth{0pt}
  \renewcommand\footrulewidth{0pt}
  \fancyheadoffset[R]{\aceRunOut}
  \fancyhead[R]{\aceRunBlock}
  \fancyfoot[C]{\aceFolio}
}

% Force the 'plain' style, used by TOC and chapter starts, to match 'mystyle'
\fancypagestyle{plain}{
  \fancyhf{}
  \renewcommand\headrulewidth{0pt}
  \renewcommand\footrulewidth{0pt}
  \fancyheadoffset[R]{\aceRunOut}
  \fancyhead[R]{\aceRunBlock}
  \fancyfoot[C]{\aceFolio}
}

\pagestyle{mystyle}


% ---------------------------------------------------------
% CHAPTER & SECTION FONTS
% ---------------------------------------------------------
\chapterfont{\color{dark}\fontsize{20}{16.8}\selectfont}

\setcounter{secnumdepth}{2}
\renewcommand{\thesection}{\arabic{section}}
\renewcommand{\thesubsection}{\thesection.\Alph{subsection}}

% Section headings, matching the site: the number in mono in the text
% cyan, a thin vertical divider, the title in bold sans, and a rule underneath
% made of a short cyan segment and a grey hairline. Subsections carry a
% mono number in the text cyan.
\definecolor{acecyantext}{HTML}{0B7F97}
\colorlet{acecyantint}{acecyan!22!white}
\newcommand{\aceSecNum}[1]{{\ttfamily\bfseries\color{acecyantext}#1}%
  \hspace{0.55em}{\color{acegrey!45}\rule[-0.22em]{1pt}{1.15em}}}
\newcommand{\aceSecRule}{\par\nobreak\vspace{-1.4ex}\noindent%
  {\color{acecyan}\rule[0.2pt]{2.4em}{1.4pt}}%
  {\color{acegrey!45}\rule[0.7pt]{\dimexpr\textwidth-2.4em\relax}{0.4pt}}\par\nobreak}

\titleformat{\section}
  {\sffamily\Large\bfseries}
  {\aceSecNum{\thesection}}
  {0.7em}
  {}
  [\aceSecRule]
\titlespacing*{\section}{0pt}{3.2ex plus 1ex minus .2ex}{1.8ex plus .2ex}

\titleformat{\subsection}
  {\sffamily\large\bfseries}
  {{\ttfamily\bfseries\color{acecyantext}\thesubsection}}
  {0.7em}
  {}

% ---------------------------------------------------------
% FRONT MATTER PAGE (title, authors, abstract, contents)
% The page after the cover. The title block repeats the title, subtitle and
% authors but not the dates, which the cover already carries. Each part is
% introduced by the same mono caps key used on the cover.
% ---------------------------------------------------------
% An exact vertical gap. Both of LaTeX's automatic skips are taken out of
% the way first: \parskip, which it adds when the next paragraph starts,
% and the interline glue, which otherwise stretches the gap out to a full
% baseline. What is asked for here is what lands on the page.
\newcommand{\aceFrontGap}[1]{\par\nobreak\nointerlineskip\vskip #1\relax}

% A section key with its rule tucked directly underneath, as on the cover.
\newcommand{\aceFrontKey}[1]{%
  \par\begingroup\parskip=0pt\parindent=0pt
  \noindent{\aceFacetKeyFont\color{acegrey}#1}%
  \aceFrontGap{0.13cm}%
  \noindent{\color{black!25}\rule{\textwidth}{0.4pt}}\par\nobreak
  \endgroup
  \nointerlineskip}

\makeatletter
\renewcommand{\maketitle}{%
  \thispagestyle{plain}%
  \begingroup\parskip=0pt\parindent=0pt
  \noindent{\fontsize{19}{23}\selectfont\bfseries\@title\par}%
  \@ifundefined{@subtitle}{}{%
    \ifx\@subtitle\@empty\else
      \aceFrontGap{0.20cm}%
      \noindent{\fontsize{12}{15}\selectfont\color{acegrey}\@subtitle\par}%
    \fi}%
  \ifx\aceAuthorsLine\@empty\else
    \aceFrontGap{0.62cm}%
    \aceFrontKey{AUTHORS}%
    \aceFrontGap{0.26cm}%
    \noindent{\fontsize{11}{14}\selectfont\bfseries\aceAuthorsLine\par}%
    \ifx\aceAffiliationsLine\@empty\else
      \aceFrontGap{0.20cm}%
      \noindent{\fontsize{8.5}{11}\selectfont\color{acegrey}\aceAffiliationsLine\par}%
    \fi
  \fi
  \endgroup
  \aceFrontGap{0.62cm}%
}

% The abstract, styled in place: Quarto emits it right after \maketitle.
\newcommand{\aceAbstractStart}{%
  \aceFrontKey{ABSTRACT}%
  \fontsize{9.5}{12.6}\selectfont
  \aceFrontGap{0.26cm}%
  \noindent\ignorespaces}
\newcommand{\aceAbstractEnd}{\aceFrontGap{0.62cm}}
\AtBeginDocument{%
  \@ifundefined{abstract}%
    {\newenvironment{abstract}{\aceAbstractStart}{\aceAbstractEnd}}%
    {\renewenvironment{abstract}{\aceAbstractStart}{\aceAbstractEnd}}%
}

% ---------------------------------------------------------
% CONTENTS THAT ALWAYS FIT THE PAGE
% The contents list is typeset into a box first and measured against the
% space left on the page. The largest setting that fits wins; the last
% resort is two columns, which holds roughly forty entries. Nothing here
% changes what is listed, only how tightly it is set.
% ---------------------------------------------------------
\usepackage{multicol}
\newsavebox{\aceTocBox}
\newlength{\aceTocSpace}
\newlength{\aceTocNeed}

% The entries themselves, read back from the .toc written on the last pass.
\newcommand{\aceTocEntries}{\@input{\jobname.toc}}

% Last resort for a very long report: list the numbered sections only.
\newcommand{\aceTocSectionsOnly}{%
  \let\l@subsection\@gobbletwo
  \let\l@subsubsection\@gobbletwo}

% Measure the list at a given type size and leading factor.
% #1 size, #2 leading, #3 extra setup (e.g. dropping the subsections)
\newcommand{\aceTocMeasure}[3]{%
  \sbox{\aceTocBox}{%
    \begin{minipage}{\textwidth}%
      \fontsize{#1}{#2}\selectfont #3%
      \aceTocEntries
    \end{minipage}}%
  \setlength{\aceTocNeed}{\dimexpr\ht\aceTocBox+\dp\aceTocBox\relax}}

\newcommand{\aceTocMeasureTwoCol}[3]{%
  \sbox{\aceTocBox}{%
    \begin{minipage}{\textwidth}%
      \fontsize{#1}{#2}\selectfont #3%
      \setlength{\columnsep}{1.2em}%
      \begin{multicols}{2}\aceTocEntries\end{multicols}%
    \end{minipage}}%
  \setlength{\aceTocNeed}{\dimexpr\ht\aceTocBox+\dp\aceTocBox\relax}}

% \@starttoc also opens the stream the next run reads back, so that has to
% happen even though the entries are placed by hand.
\newcommand{\aceTocOpenWrite}{%
  \if@filesw
    \expandafter\newwrite\csname tf@toc\endcsname
    \immediate\openout \csname tf@toc\endcsname \jobname.toc\relax
  \fi}

\renewcommand{\tableofcontents}{%
  \par
  % Space left on this page, guarding against being called before the page
  % has any material on it.
  \setlength{\aceTocSpace}{\dimexpr\pagegoal-\pagetotal\relax}%
  \ifdim\pagegoal=\maxdimen
    \setlength{\aceTocSpace}{\dimexpr\textheight-\pagetotal\relax}%
  \fi
  \aceFrontKey{SECTIONS}%
  \aceFrontGap{0.26cm}%
  % what the key, its rule and the bottom margin have already taken
  \addtolength{\aceTocSpace}{-1.3cm}%
  % Largest setting that fits, in order: full size, tighter, two columns,
  % and finally the numbered sections on their own.
  \aceTocMeasure{10}{15}{}%
  \ifdim\aceTocNeed>\aceTocSpace \aceTocMeasure{9}{13.2}{}\fi
  \ifdim\aceTocNeed>\aceTocSpace \aceTocMeasure{8}{11.6}{}\fi
  \ifdim\aceTocNeed>\aceTocSpace \aceTocMeasureTwoCol{8.5}{12}{}\fi
  \ifdim\aceTocNeed>\aceTocSpace \aceTocMeasureTwoCol{7.5}{10.5}{}\fi
  \ifdim\aceTocNeed>\aceTocSpace \aceTocMeasure{9.5}{14}{\aceTocSectionsOnly}\fi
  \ifdim\aceTocNeed>\aceTocSpace \aceTocMeasure{8}{11.6}{\aceTocSectionsOnly}\fi
  \ifdim\aceTocNeed>\aceTocSpace \aceTocMeasureTwoCol{8}{11.6}{\aceTocSectionsOnly}\fi
  \noindent\usebox{\aceTocBox}\par
  \aceTocOpenWrite
}
\makeatother

% ---------------------------------------------------------
% BUILD THE COVER PAGE
% ---------------------------------------------------------
\makeatletter
\AtBeginDocument{
    \thispagestyle{empty} % Ensure no headers/footers on the cover

    % Title and subtitle are set into a box, and their width measured, before
    % the picture starts: TeX cannot measure anything inside a tikzpicture.
    % The band behind them then takes its height from the type rather than
    % from a guess, so a title that wraps to two lines makes the band taller
    % and pushes the dates down with it.
    \sbox{\aceTitleBox}{%
        \begin{minipage}{\aceTitleMeasure}\raggedright
        {\fontsize{28}{32}\selectfont\bfseries\color{acecoverink}\@title\par}%
        \vskip 0.24cm
        {\fontsize{15}{18}\selectfont\color{acecoversub}\@subtitle\par}%
        \end{minipage}}%
    \setlength{\aceTitleHeight}{\dimexpr\ht\aceTitleBox+\dp\aceTitleBox\relax}%

    % How far the band runs: the longer of the two lines, plus the margin it
    % starts from and a little air at the end. A title that fills the measure
    % stops the band at the measure.
    \settowidth{\aceTitleWidth}{\fontsize{28}{32}\selectfont\bfseries\@title}%
    \settowidth{\aceTitleRun}{\fontsize{15}{18}\selectfont\@subtitle}%
    \ifdim\aceTitleRun>\aceTitleWidth \setlength{\aceTitleWidth}{\aceTitleRun}\fi
    \ifdim\aceTitleWidth>\aceTitleMeasure \setlength{\aceTitleWidth}{\aceTitleMeasure}\fi
    \addtolength{\aceTitleWidth}{2.9cm}% 2cm margin plus 0.9cm of air
    \setlength{\aceBandTop}{-8.70cm}%
    \setlength{\aceBandBottom}{\dimexpr -9.2cm - \aceTitleHeight - 0.50cm\relax}%

    % Links are lighter on the dark cover than in the body; restored below.
    \hypersetup{urlcolor=acecoverlink}

    \begin{tikzpicture}[remember picture,overlay]

        % The cover is a dark page, edge to edge. The rectangle is bled 3mm
        % past every edge: a fill that stops exactly on the page boundary
        % leaves viewers a hairline of white where the edge pixel is blended
        % with the page background. Anything outside the page is clipped.
        \fill[acecoverbg] ([xshift=-3mm, yshift=-3mm]current page.south west)
            rectangle ([xshift=3mm, yshift=3mm]current page.north east);
        
        % --- Masthead -------------------------------------------------------
        % The cover is a dark page: logo top left, document label top right,
        % then the title block on a band that bleeds off the left edge. The
        % glitch band that used to fill the middle was removed; scripts/
        % gen_cover_band.py still generates it if it is ever wanted back.
        \node at ([xshift=2cm, yshift=-3.1cm]current page.north west) [anchor=west]
            {\includegraphics[width=5.6cm]{brand/logo_blackbg_text.png}};

        % Document label, balancing the logo across the page. The rule is cut
        % to the width of the label itself.
        \node (aceCoverLabel) at ([xshift=-2cm, yshift=-2.95cm]current page.north east)
            [anchor=south east, text=acecoverink] {\aceCoverLabelFont AUDIT REPORT};
        \draw[acecyan, line width=1.4pt] (aceCoverLabel.south west) -- (aceCoverLabel.south east);
        \node at ([xshift=-2cm, yshift=-3.48cm]current page.north east)
            [anchor=north east, text=acecovermute] {\aceCoverIDFont \aceID};

        % --- ECG trace ------------------------------------------------------
        % A rhythm strip sitting just above the title band, bleeding off the
        % left edge with it and stopping at the same right edge. Drawn in a flat
        % grey a step above the band's own value, rather than in white at low
        % opacity: transparency made every overlap brighter than the line. The
        % clip keeps the trace out of the band however wide the title runs.
        \begin{scope}
            \clip ([xshift=-3mm, yshift=\dimexpr\aceBandTop+1.70cm\relax]current page.north west)
               rectangle ([xshift=\aceTitleWidth, yshift=\aceBandTop]current page.north west);
            \draw[aceecg, line width=1.4pt, line cap=round, line join=round]
              ([xshift=-0.3cm, yshift=\dimexpr\aceBandTop+0.55cm\relax]current page.north west)
              \aceEcgStrip;
        \end{scope}

        % --- Title block ----------------------------------------------------
        % The band bleeds off the left edge and stops just past the longer of
        % the two lines; its height and width were measured above.
        \fill[acetitleband]
            ([xshift=-3mm, yshift=\aceBandTop]current page.north west) rectangle
            ([xshift=\aceTitleWidth, yshift=\aceBandBottom]current page.north west);

        \node at ([xshift=2cm, yshift=-9.2cm]current page.north west)
            [anchor=north west, inner sep=0pt] {\usebox{\aceTitleBox}};

        % Dates, a fixed distance below the band whatever its height.
        \node at ([xshift=2cm, yshift=\dimexpr\aceBandBottom-0.62cm\relax]current page.north west)
            [anchor=north west, inner sep=0pt, text=acecoverink] {%
            \begin{tabular}{@{}l@{\hspace{1.5cm}}l@{}}
            {\aceFacetKeyFont\color{acefacetpopink}PUBLISHED} & {\aceFacetKeyFont\color{acefacetdesink}MODIFIED}\\[0.16cm]
            {\aceCoverDateFont\@date} & {\aceCoverDateFont\dateModified}
            \end{tabular}};

        \coverAuthorsBlock

        % Facet badges, left side, generated by fields.lua from the report's
        % population/treatment/outcome/design fields. Centered in the same
        % white space as the author block.
        \coverCategoriesBlock
        
        % --- Document information -------------------------------------------
        % A key column and a value column between two hairlines, with the
        % disclaimer running as prose underneath. The report ID is not
        % repeated here: it is set in the corner with the document label.
        %
        % The block is anchored at its bottom rule (-12.71cm) and rows are
        % stacked upward, so it stays contiguous whichever DOI fields a
        % report has. \aceInfoY ends one row above the topmost row, which is
        % what places the top rule. fields.lua mirrors this arithmetic to
        % know where the author column has to stop.
        \setlength{\aceInfoY}{-12.05cm}
        \node at ([xshift=2cm, yshift=\aceInfoY]current page.west) [anchor=north west, inner sep=0pt, text=acecovermute] {\aceFacetKeyFont LICENSE};
        \node at ([xshift=6.6cm, yshift=\aceInfoY]current page.west) [anchor=north west, text=acecoverink] {\aceInfoValueFont Creative Commons Attribution 4.0 International (CC BY 4.0)};
        \addtolength{\aceInfoY}{0.62cm}
        \node at ([xshift=2cm, yshift=\aceInfoY]current page.west) [anchor=north west, inner sep=0pt, text=acecovermute] {\aceFacetKeyFont COPYRIGHT};
        \node at ([xshift=6.6cm, yshift=\aceInfoY]current page.west) [anchor=north west, text=acecoverink] {\aceInfoValueFont © 2026 Matthew B. Jané};
        \addtolength{\aceInfoY}{0.62cm}
        \ifx\aceConceptDOI\empty\else
        \node at ([xshift=2cm, yshift=\aceInfoY]current page.west) [anchor=north west, inner sep=0pt, text=acecovermute] {\aceFacetKeyFont ALL VERSIONS};
        \node at ([xshift=6.6cm, yshift=\aceInfoY]current page.west) [anchor=north west] {\aceInfoMonoFont\href{https://doi.org/\aceConceptDOI}{https://doi.org/\aceConceptDOI}};
        \addtolength{\aceInfoY}{0.62cm}
        \fi
        \ifx\aceDOI\empty\else
        \node at ([xshift=2cm, yshift=\aceInfoY]current page.west) [anchor=north west, inner sep=0pt, text=acecovermute] {\aceFacetKeyFont VERSION DOI};
        \node at ([xshift=6.6cm, yshift=\aceInfoY]current page.west) [anchor=north west] {\aceInfoMonoFont\href{https://doi.org/\aceDOI}{https://doi.org/\aceDOI}};
        \addtolength{\aceInfoY}{0.62cm}
        \fi

        \setlength{\aceInfoTop}{\aceInfoY}\addtolength{\aceInfoTop}{-0.14cm}
        \draw[acecovermute!55, line width=0.5pt] ([xshift=2cm, yshift=\aceInfoTop]current page.west) -- ([xshift=-2cm, yshift=\aceInfoTop]current page.east);
        \draw[acecovermute!55, line width=0.5pt] ([xshift=2cm, yshift=-12.71cm]current page.west) -- ([xshift=-2cm, yshift=-12.71cm]current page.east);
        \node at ([xshift=2cm, yshift=-12.91cm]current page.west) [anchor=north west, text width=16.6cm, text=acecovermute]
            {\aceInfoNoteFont The views expressed herein are solely those of the authors and do not represent their affiliated organizations. Full disclaimers and editorial policy at \href{https://auditclinicalevidence.org/disclaimers.html}{auditclinicalevidence.org/disclaimers}.};

    \end{tikzpicture}

    \hypersetup{urlcolor=highlight}

    % End the cover page
    \clearpage
    
    % Force the NEXT page to be labeled as Page 1
    \setcounter{page}{1}
}
\makeatother

% ---------------------------------------------------------
% CUSTOM FONTS
% ---------------------------------------------------------
\setmainfont{Roboto}[
    Path=../_extensions/ace-formatting/acePDF/Roboto/static/,
    Scale=0.9,
    Extension = .ttf,
    UprightFont=*-Regular,
    BoldFont=*-Bold,
    ItalicFont=*-Italic,
    ]
    
\setsansfont{Roboto}[
    Path=../_extensions/ace-formatting/acePDF/Roboto/static/,
    Scale=0.9,
    Extension = .ttf,
    UprightFont=*-Regular,
    BoldFont=*-Bold,
    ItalicFont=*-Italic,
    ]

\setmonofont{RobotoMono}[
    Path=../_extensions/ace-formatting/acePDF/Roboto_Mono/static/,
    Scale=0.9,
    Extension = .ttf,
    UprightFont=*-Medium,
    BoldFont=*-Bold,
    ItalicFont=*-Italic,
    ]
    
\setmathfont{Concrete-Math.otf}[
  % Concrete Math: Knuth's low-contrast Concrete, the same math font as the
  % website (styles.css, "MATH"). Even strokes that sit at Roboto's weight.
  % Bundled in the extension's fonts/ folder, like Roboto - no tlmgr needed.
  Path = ../_extensions/ace-formatting/acePDF/fonts/,
  Scale = MatchLowercase
]

% Calligraphic capitals come from Euler Math, Concrete's original partner: Zapf's Euler
% Script is upright-leaning with a uniform stroke, so it stays legible and
% sturdy beside a sans body where Computer Modern's hairlines look wispy.
% Alternatives, same line: STIXTwoMath-Regular.otf (formal roundhand),
% latinmodern-math.otf (Computer Modern).
\setmathfont{Euler-Math.otf}[
  range = {cal, bfcal},
  Path = ../_extensions/ace-formatting/acePDF/fonts/,
  Scale = MatchLowercase
]

% A float that gets a page to itself (a large table or figure) sits at the
% top of that page, not centered vertically with white space above it.
\makeatletter
\setlength{\@fptop}{0pt}
\setlength{\@fpsep}{12pt plus 2fil}
\setlength{\@fpbot}{0pt plus 1fil}
\makeatother
